\chapter{Workload: Pod, Deployment, StatefulSet}
\label{ch:workloads}

\section{Pod}

Pod là đơn vị mà Kubernetes lập lịch: một hoặc nhiều container dùng
chung network namespace (một IP, \code{localhost} giữa chúng) và volume.
Hầu như mọi pod ở đây chỉ có một container; các mẫu nhiều container
(sidecar làm proxy hay đẩy log) mới quan trọng khi bạn gặp service mesh.
Pod \emph{sinh ra để chết} --- chúng không bao giờ được chữa, chỉ được
thay, và IP của chúng chết theo. Mọi thứ còn lại trong Phần~V ---
controller, Service, PVC --- tồn tại để xây hành vi bền vững từ những
pod dùng xong bỏ.

\section{Deployment và rollout}

Bạn không bao giờ tạo pod trần. Một \emph{Deployment} khai báo ``N
replica của mẫu pod này''; nó quản lý một \emph{ReplicaSet} (bộ đếm),
và ReplicaSet quản lý các pod. Ba tầng gián tiếp là thứ làm cho việc
cập nhật trở nên thanh lịch: đổi mẫu (một tag image mới), Deployment
tạo một ReplicaSet \emph{mới}, tăng dần nó lên trong khi cái cũ giảm
xuống:

\begin{figure}[h]
\centering
\begin{tikzpicture}[node distance=3.5mm and 9mm]
  \node[box] (d) {Deployment\\\scriptsize course-service};
  \node[box, below left=of d] (rs1) {ReplicaSet (cũ)\\\scriptsize image :a1b2c3d};
  \node[box, below right=of d] (rs2) {ReplicaSet (mới)\\\scriptsize image :e4f5a6b};
  \node[boxfill, below=of rs1] (p1) {pod \scriptsize(đang tắt)};
  \node[boxfill, below=of rs2] (p2) {pod \scriptsize(khởi động)};
  \draw[flow] (d) -- (rs1); \draw[flow] (d) -- (rs2);
  \draw[flowdash] (rs1) -- (p1); \draw[flow] (rs2) -- (p2);
\end{tikzpicture}
\caption{Một rolling update đang giữa chừng: hai ReplicaSet, lưu lượng đang dịch chuyển.}
\end{figure}

Rollout chỉ tiến tới khi pod mới trở nên \emph{Ready} (probe của
Chương~17 --- đây là nơi chúng chứng tỏ giá trị), nên một image hỏng
làm rollout dừng lại thay vì kéo sập service; và
\code{kubectl rollout undo} trỏ ngược về ReplicaSet trước, thứ được giữ
lại chính là vì mục đích này. \code{kubectl rollout status} của
pipeline (Chương~23) là CI đang chờ bộ máy này kết thúc.

Với một replica --- cụm này --- rolling update vẫn nghĩa là
khởi-động-mới-rồi-tắt-cũ, một khoảng chồng lấn hai pod ngắn ngủi tốn
một thoáng gấp đôi bộ nhớ. \code{maxUnavailable}/\code{maxSurge} tinh
chỉnh điệu nhảy đó ở quy mô lớn.

\section{StatefulSet: khi danh tính quan trọng}

Với cơ sở dữ liệu, tính hoán đổi được là hoàn toàn sai: pod nào lên
cũng phải là \emph{đúng Postgres đó} với \emph{đúng dữ liệu đó}.
StatefulSet cung cấp thứ mà Deployment từ chối: tên ổn định
(\code{postgres-course-0}, không phải hậu tố ngẫu nhiên), PVC riêng cho
từng replica qua \code{volumeClaimTemplates} (Chương~19) sống lâu hơn
pod và gắn lại vào pod kế nhiệm, và thay thế có thứ tự, từng cái một.
Từ \code{postgres-course.yaml}:

\begin{lstlisting}[language=yaml]
kind: StatefulSet
spec:
  serviceName: postgres-course
  replicas: 1
  volumeClaimTemplates:          # a PVC per replica, not per manifest
    - metadata: { name: data }
      spec:
        accessModes: [ReadWriteOnce]
        resources: { requests: { storage: 5Gi } }
\end{lstlisting}

Quy tắc quyết định: \emph{tiến trình} có giữ trạng thái mà tiến trình
kế tiếp phải thừa kế không? Cơ sở dữ liệu thì có; Spring service thì
không --- trạng thái của chúng nằm trong các cơ sở dữ liệu kia, điều
làm chúng thành Deployment, và điều khiến chúng mở rộng và rollout được.
(Tính phi trạng thái của JWT ở Chương~5 chính là thuộc tính này ở tầng
trên.)

\begin{pitfall}[mất dữ liệu kiểu cơ-sở-dữ-liệu-trong-Deployment]
Triệu chứng: ai đó chạy cơ sở dữ liệu như một Deployment với volume
thường; chạy tốt hàng tháng; rồi một sự cố node tạo ra replica
\emph{thứ hai} trỏ vào cùng bộ file, và cơ sở dữ liệu hỏng. Nguyên
nhân: hợp đồng của Deployment là ``các replica hoán đổi được'' và
rolling update của nó chạy cũ và mới \emph{đồng thời} trong chốc lát
--- chí mạng khi cả hai cùng mở một thư mục dữ liệu. Thứ tự của
StatefulSet + claim riêng từng replica tồn tại để điều này bất khả. Bài
học khái quát: chọn controller theo \emph{hợp đồng} của nó, không phải
theo file YAML nào bạn thấy trước.
\end{pitfall}

\section{Các loại workload khác}

Cho đầy đủ, các controller còn lại, chưa cái nào được dùng ở đây:
\textbf{DaemonSet} (một pod mỗi node --- bộ gom log, metric của node),
\textbf{Job} (chạy tới khi xong --- Flyway-như-bước-CI của Chương~7 sẽ
là một Job), \textbf{CronJob} (Job theo lịch --- sao lưu, mảnh ghép còn
thiếu của Chương~19).

\begin{handson}[xem một rollout diễn ra]
\begin{lstlisting}[language=bash]
kubectl -n ts get rs -l app=course-service          # one ReplicaSet
kubectl -n ts set image deployment/course-service \
    course-service=localhost:5000/course-service:1.0.0-SNAPSHOT
kubectl -n ts rollout status deployment/course-service
kubectl -n ts get rs -l app=course-service          # two -- old kept at 0
kubectl -n ts rollout undo deployment/course-service
\end{lstlisting}
Lệnh \code{get rs} thứ hai là điểm mấu chốt: ReplicaSet cũ vẫn tồn tại,
scale về không --- đích của undo, mặc định được giữ cho mười bản sửa
đổi gần nhất.
\end{handson}

\begin{quiz}
\begin{enumerate}
  \item Vì sao Deployment quản lý ReplicaSet thay vì quản lý pod trực
  tiếp? Tính năng nào rơi ra từ tầng gián tiếp đó?
  \item Phát biểu khác biệt về hợp đồng giữa Deployment và StatefulSet,
  mỗi cái một câu, rồi phân loại: Eureka server, một cache Redis, một
  Kafka broker, api-gateway.
  \item Trong rolling update, điều kiện nào quyết định việc tắt pod cũ,
  và bộ máy của chương nào cung cấp điều kiện đó?
  \item Giải thích chính xác cơ chế khiến chạy Postgres dưới Deployment
  rước lấy hỏng dữ liệu trong khi StatefulSet thì không.
\end{enumerate}
\end{quiz}
